\documentclass{aastex701}

\begin{document}

\title{A Giant Transiting Planet Candidate on a 3.96 Day Orbit around TIC 28812587}

\author[0000-0000-0000-0000]{Carl Drinkwater}
\affiliation{Independent researcher, Shropshire, UK}
\email{carl@29degrees.co.uk}

\begin{abstract}
I report a transit signal on TIC 28812587 (2MASS J16291561$-$3703031, $T=11.07$), a K star of 4850 K at 200 pc.
I found it in a box least squares search of the Quick-Look Pipeline light curves of TESS Sector 103, and recovered it
in the TESS-SPOC light curves of Sectors 12 and 39. A joint fit gives a period of 3.9634883 d, a depth of 5807 ppm and
a radius of $11.2 \pm 1.1\,R_\oplus$. The Sector 103 difference image places the transit on the target. The light
curves show no secondary eclipse and no odd-even depth difference. TRICERATOPS gives a false positive probability of
$0.11 \pm 0.02$. Gaia DR3 flags the star as radial velocity variable, which a planet cannot cause, so the system needs
radial velocity follow-up and high-resolution imaging.
\end{abstract}

\section{Detection and vetting}

I searched the Quick-Look Pipeline \citep[QLP;][]{huang2020a,huang2020b,kunimoto2021} light curves of TESS
\citep{ricker2015} Sector 103 for stars brighter than $T=11.5$. I removed trends with a biweight filter of 0.5 d
\citep{hippke2019} and searched 0.5 to 14 d with box least squares \citep{kovacs2002}. TIC 28812587 \citep{stassun2019}
shows a signal at 3.9628 d with a multiple event statistic of 13.5. It passes the flux-level false alarm and false
positive tests of LEO-vetter \citep{kunimoto2025}, and matches no TOI \citep{guerrero2021}, CTOI, SPOC threshold
crossing event or TESS eclipsing binary \citep{prsa2022} in the lists of 2026 October 26.

The TESS-SPOC light curves \citep{caldwell2020} of Sectors 12 and 39 show the same transit at signal-to-noise ratios of
23.5 and 28.4. TESS observed the star in no other released sector. A joint box least squares search of the three
sectors finds the signal at 3.96349 d with a signal-to-noise ratio of 46.9, against 15.4 for the highest of the null
trials.

\section{Transit fit}

I fitted a transit model \citep{kreidberg2015} to the three sectors with emcee \citep{foremanmackey2013}, with the
cadence of each sector. Figure~\ref{fig:transit} shows the fit. The fit gives:
period $3.9634883 \pm 0.0000025$ d;
mid-transit time $2461156.1382 \pm 0.0012$ BJD$_\mathrm{TDB}$;
depth $5807 \pm 187$ ppm;
duration $2.056 \pm 0.098$ h;
$R_p/R_\star = 0.074 \pm 0.005$;
impact parameter $0.71 \pm 0.24$;
$a/R_\star = 12.0 \pm 2.8$.
With the TIC stellar radius of $1.39 \pm 0.10\,R_\odot$, the companion radius is $11.2 \pm 1.1\,R_\oplus$
($1.00\,R_\mathrm{J}$). The stellar density from the fit is 5.0 times the TIC value. A smaller star than the TIC
gives would explain this, and would make the companion smaller.

The TIC contamination ratio is 0.58, and the depths of the pipeline light curves differ: $4955 \pm 261$, $5694 \pm 237$
and $7203 \pm 300$ ppm in Sectors 103, 12 and 39. Aperture photometry with the same apertures on the TESScut cutouts
\citep{brasseur2019} gives 2921 to 4599 ppm, closer between sectors ($p = 0.023$). I attribute most of the
difference to the dilution corrections of the two pipelines.

\section{False positive tests}

The odd and even transits differ by 0.45$\sigma$. The flux at phase 0.5 is $84 \pm 131$ ppm. The difference image
\citep{bryson2013} of Sector 103 places the transit on the target at a signal-to-noise ratio of 16.1. The images of
Sectors 12 and 39 are too faint to locate it. Two Gaia DR3 stars could cause the dip with eclipses of 3\% and 4\%: one
of $G=14.0$ at $34''$, and one of $G=14.3$ at $14''$ that VSX \citep{watson2006} lists as a semi-regular variable.
TRICERATOPS \citep{giacalone2021} gives a false positive probability of $0.11 \pm 0.02$ and a nearby false positive
probability of 0 over five runs, with a transiting planet on the target as the most likely scenario.

Gaia DR3 \citep{gaia2023} gives a RUWE of 0.98 and a non\_single\_star flag of 0. Its radial velocities vary,
with a $\chi^2$ p-value of 0 at the precision of the catalogue over 14 epochs and an epoch scatter of 5.5 km\,s$^{-1}$ \citep{katz2023}. A planet would give
about 0.1 km\,s$^{-1}$, which Gaia cannot detect. If the transiting body causes the variation, the scatter limits its
mass to below 0.051 $M_\odot$, so it is at most a brown dwarf. A third body, or noise near the faint limit of the
Gaia spectrometer, can also cause it. VSX lists the target as a rotational variable.

\section{Conclusion}

TIC 28812587 shows a transit of a $11.2 \pm 1.1\,R_\oplus$ companion on a 3.96 d orbit. The companion may be a giant
planet or a brown dwarf. The Gaia radial velocity variation, the high stellar density from the fit and the
contamination of the aperture leave open a bound or blended stellar companion. Radial velocities can measure the mass
and test for a third body, and high-resolution imaging can test the neighbours.

\begin{figure}[ht!]
\plotone{figure_1.pdf}
\caption{Transit of TIC 28812587 in the Sector 103 QLP (left), Sector 12 and Sector 39 TESS-SPOC (middle) light
curves, with the joint fit (red) and 11-minute bins (black). Right: the flux at phase 0.5, with the transit duration
shaded.
\label{fig:transit}}
\end{figure}

\begin{acknowledgments}
This note uses data from the TESS mission, obtained from MAST. Claude (Anthropic) wrote the analysis code and a draft of
this note under my direction. The code and results are at https://github.com/carldr/tess-transiting-planet-search.
\end{acknowledgments}

\facilities{TESS, Gaia}

\software{astropy, astroquery, batman \citep{kreidberg2015}, emcee \citep{foremanmackey2013}, LEO-vetter
\citep{kunimoto2025}, lightkurve, TRICERATOPS \citep{giacalone2021}, wotan \citep{hippke2019}}

\bibliography{references}{}
\bibliographystyle{aasjournalv7}

\end{document}
